\section{Discussion}
\label{sec:discussion}

\subsection{When Does Persistent Memory Help?}

Cross-session memory is most useful for tasks where the relevant
knowledge is (a) acquired during a previous episode, (b) stable
over time (e.g., object locations, user preferences), and (c) not
trivially re-observable. In our analysis of existing memory
benchmarks (Section~\ref{sec:related}), none of RMBench,
MemoryBench, MIKASA, or RoboCerebra contains genuine cross-session
tasks: each task starts from a fresh memory state. We view this as a
structural gap in evaluation that MNEMO-Bench is designed to close
through its \emph{cross-session} axis.

\subsection{Interpretability as a First-Class Feature}

A recurrent theme in our design is that memory should be
\emph{readable}. The semi-structured language format makes failure
diagnosis straightforward: a developer can dump the cold memory at
any point during an episode and read it as English. This is a
stark contrast to dense-embedding memories like MemoryVLA's PCMB
or ReMem-VLA's recurrent queries, where post-hoc interpretation is
restricted to attention-heatmap visualization. We argue that
interpretability is not just a debugging convenience but a
prerequisite for deployment in safety-critical domains.

\subsection{Limitations}

\paragraph{No hot-memory experiments yet.}
The current release does not include a hot-memory encoder;
integration with Mem-0's sliding window suffices for our prompt-level
experiments. Implementing a MEM-style video encoder with tactile
fusion is scheduled for v0.3 of the codebase and will enable
full-policy training experiments.

\paragraph{Embedder choice matters.}
The toy \texttt{HashingEmbedder} that ships with v0.2 achieves a
surprisingly strong Recall@1=75\% on the synthetic benchmark, but
semantic generalization to paraphrases with distinct lexical forms
requires a real encoder. Production deployments should use the
\texttt{SentenceTransformerEmbedder} we ship or the text encoder of
the underlying VLA. We document the substitution clearly.

\paragraph{Rule-based summarization ceiling.}
Our initial episode summarizer uses deterministic rules to extract
narratives and facts. It is adequate for the interface test-bed but
lacks the compression and abstraction an LLM-based summarizer would
provide. v0.2 will add a batched LLM summarizer that maintains the
deterministic fallback for offline or privacy-sensitive settings.

\paragraph{Benchmark coverage.}
MNEMO-Bench's v1.0 draft selects representative tasks from each of
the four source benchmarks, but the exact selection remains
calibrated empirically. We will release the task selection criteria
and an open leaderboard with the camera-ready.

\subsection{Ethical Considerations}

Persistent memory raises unique safety concerns: a robot that
remembers user interactions across sessions must handle
privacy-sensitive information responsibly. Our schema captures
facts with explicit provenance and confidence, enabling
session-scoped redaction and audit logging. We document a
redaction toolkit in the appendix and recommend opt-in memory for
consumer deployments.
